Use the skymap ``Map-OM02'' to answer the questions below.

The constellation lines and boundaries (as per IAU standards) of two
constellations denoted as C1 and C2 are shown in the full-sky map. Alternative
depictions of the same constellations according to a few cultures are also
shown on the right panel for your reference, if needed. A certain coordinate
grid is also shown.

% FIGURE NEEDED: full-sky map "Map-OM02" with constellations C1, C2, a
% coordinate grid, red squares and blue circles, and shaded constellation
% areas (2025_SKYMAP paper)

\begin{description}
\item[(OM02.1)] Identify the constellations C1 and C2 and write their names
  (of Latin origin) or IAU abbreviations in the table in the Summary
  Answersheet.
\item[(OM02.2)] Three empty red squares and three empty blue circles are shown
  on the map. There is a grid line passing through each of these squares and
  circles.
  \begin{description}
  \item[(OM02.2a)] The lines that pass through the red squares are lines of
    constant Ecliptic latitude ($\beta$) / Ecliptic longitude ($\lambda$) /
    Declination ($\delta$) / Right Ascension ($\alpha$) / Galactic latitude
    ($b$) / Galactic longitude ($l$). Tick ($\checkmark$) the correct option
    in the Summary Answersheet.
  \item[(OM02.2b)] The lines that pass through the blue circles are lines of
    constant Ecliptic latitude ($\beta$) / Ecliptic longitude ($\lambda$) /
    Declination ($\delta$) / Right Ascension ($\alpha$) / Galactic latitude
    ($b$) / Galactic longitude ($l$). Tick ($\checkmark$) the correct option
    in the Summary Answersheet.
  \end{description}
\item[(OM02.3)] Identify the North and South poles of the grid. Label these
  points as ``N'' and ``S'', respectively, on the map ``Map-OM02''.
\item[(OM02.4)] Two of the following are present in the given sky map.
  Identify these by marking with appropriate symbols (shown below) on the
  corresponding entire curve/line.
  \begin{description}
  \item[(i)] Ecliptic (small bars)
  \item[(ii)] Celestial Equator (small circles)
  \item[(iii)] Galactic Equator (small crosses)
  \end{description}
\item[(OM02.5)] Mark the Vernal Equinox (VE) and Autumnal Equinox (AE) on the
  grid with $\otimes$ and write VE and AE beside them, respectively.
\item[(OM02.6)] Indicate the direction of the Sun's annual motion by drawing
  an arrow close to the Vernal Equinox.
\item[(OM02.7)] Write the values, in appropriate units, inside each red square
  and blue circle given on the ``Map-OM02'', of the corresponding grid lines
  passing through them.
\item[(OM02.8)] The location of 4 constellations (apart from C1 and C2) are
  shown on the grid by light-green shaded areas. Consider the following list
  of constellations.

  Aquarius (Aqr), Cygnus (Cyg), Leo (Leo), Orion (Ori), Perseus (Per),
  Sagittarius (Sgr).

  On the map ``Map-OM02'', label the appropriate shaded areas with the IAU
  abbreviations of the constellations that are present in the above list.
  Mark a cross ($\times$) on those shaded areas, if any, which do not appear
  in the above list.
\end{description}
